
\documentclass{article}
%%%%%%%%%%%%%%%%%%%%%%%%%%%%%%%%%%%%%%%%%%%%%%%%%%%%%%%%%%%%%%%%%%%%%%%%%%%%%%%%%%%%%%%%%%%%%%%%%%%%%%%%%%%%%%%%%%%%%%%%%%%%
\usepackage{amsfonts}
\usepackage{achemso}

%TCIDATA{OutputFilter=LATEX.DLL}
%TCIDATA{Version=4.10.0.2363}
%TCIDATA{Created=Wednesday, November 03, 2004 11:56:22}
%TCIDATA{LastRevised=Wednesday, November 03, 2004 11:57:00}
%TCIDATA{<META NAME="GraphicsSave" CONTENT="32">}
%TCIDATA{<META NAME="DocumentShell" CONTENT="Articles\SW\American Chemical Society">}
%TCIDATA{CSTFile=LaTeX Article (bright).cst}

\newtheorem{theorem}{Theorem}
\newtheorem{acknowledgement}[theorem]{Acknowledgement}
\newtheorem{algorithm}[theorem]{Algorithm}
\newtheorem{axiom}[theorem]{Axiom}
\newtheorem{case}[theorem]{Case}
\newtheorem{claim}[theorem]{Claim}
\newtheorem{conclusion}[theorem]{Conclusion}
\newtheorem{condition}[theorem]{Condition}
\newtheorem{conjecture}[theorem]{Conjecture}
\newtheorem{corollary}[theorem]{Corollary}
\newtheorem{criterion}[theorem]{Criterion}
\newtheorem{definition}[theorem]{Definition}
\newtheorem{example}[theorem]{Example}
\newtheorem{exercise}[theorem]{Exercise}
\newtheorem{lemma}[theorem]{Lemma}
\newtheorem{notation}[theorem]{Notation}
\newtheorem{problem}[theorem]{Problem}
\newtheorem{proposition}[theorem]{Proposition}
\newtheorem{remark}[theorem]{Remark}
\newtheorem{solution}[theorem]{Solution}
\newtheorem{summary}[theorem]{Summary}
\newenvironment{proof}[1][Proof]{\noindent\textbf{#1.} }{\ \rule{0.5em}{0.5em}}
\input{tcilatex}

\begin{document}

\title{\texttt{achemso}\\
-- A BibTeX style for Chemistry Publications\thanks{%
This document describes \textsf{achemso} v.1.0{} and was last updated on
1998/06/01 .}}
\author{Mats Dahlgren \\
%EndAName
(\texttt{matsd@homenet.se})}
\date{1998/06/01}
\maketitle

\begin{abstract}
The \textsf{achemso} package provides a BibTeX style in accordance with the
requirements of the journals of The American Chemical Society. It consists
of both a BibTeX style file and a LaTeX package file. Also provided, is a
BibTeX style file to be used for bibliography data base listings via an
option to the LaTeX package \textsf{achemso}. \newline
Package: Copyright \copyright\ 1995-1998 by Mats Dahlgren. All rights
reserved.
\end{abstract}

\section{About This Shell}

This shell document provides a sample layout of an American Chemical Society
article.

The front matter has a number of sample entries that you should replace with
your own. Replace this text with your own. You may delete all of the text in
this document to start with a blank document.

Further instructions and information about the typesetting specification
used by this shell document can be found starting at Section \ref{achemsodoc}%
, which is included in this shell document from the original documentation
file included with the typesetting specification.

Changes to the typeset format of this shell and its associated \LaTeX{}
formatting file (\texttt{article.cls }and \texttt{achemso.sty}) are not
supported by MacKichan Software, Inc. If you wish to make such changes,
please consult the \LaTeX{} manuals or a local \LaTeX{} expert.

If you modify this document and save it as ``American Chemical Society.shl''
in the \texttt{Shells\TEXTsymbol{\backslash}Articles\TEXTsymbol{\backslash}SW%
} directory, it will become your new American Chemical Society shell.

The typesetting specification selected by this shell document uses the
default class options. There are a number of class options supported by this
typesetting specification. The available options include setting the paper
size and the point size of the font used in the body of the document. Select 
\textsf{Typeset, Options and Packages}, the \textsf{Class Options} tab and
then click the \textsf{Modify} button to see the class options that are
available for this typesetting specification.

\section{Features of This Shell}

\subsection{Subsection}

Use the Section tag for major sections, and the Subsection tag for
subsections.

\subsubsection{Subsubsection}

This is just some harmless text under a subsubsection.

\paragraph{Subsubsubsection}

This is just some harmless text under a subsubsubsection.

\subparagraph{Subsubsubsubsection}

This is just some harmless text under a subsubsubsubsection.

\subsection{Tags}

You can apply the logical markup tag \emph{Emphasized}.

You can apply the visual markup tags \textbf{Bold}, \textit{Italics}, 
\textrm{Roman}, \textsf{Sans Serif}, \textsl{Slanted}, \textsc{Small Caps},
and \texttt{Typewriter}.

You can apply the special, mathematics only, tags $\mathbb{BLACKBOARD}$ $%
\mathbb{BOLD}$, $\mathcal{CALLIGRAPHIC}$, and $\mathfrak{fraktur}$. Note
that blackboard bold and calligraphic are correct only when applied to
uppercase letters A through Z.

You can apply the size tags {\tiny tiny}, {\scriptsize scriptsize}, 
{\footnotesize footnotesize}, {\small small}, {\normalsize normalsize}, 
{\large large}, {\Large Large}, {\LARGE LARGE}, {\huge huge} and {\Huge Huge}%
.

\QTP{Body Math}
This is a Body Math paragraph. Each time you press the Enter key, Scientific
WorkPlace switches to mathematics mode. This is convenient for carrying out
``scratchpad'' computations.

Following is a group of paragraphs marked as Body Quote. This environment is
appropriate for a short quotation or a sequence of short quotations.

\begin{quote}
The buck stops here. \emph{Harry Truman}

Ask not what your country can do for you; ask what you can do for your
country. \emph{John F Kennedy}

I am not a crook. \emph{Richard Nixon}

It's no exaggeration to say the undecideds could go one way or another. 
\emph{George Bush}

I did not have sexual relations with that woman, Miss Lewinsky. \emph{Bill
Clinton}
\end{quote}

The Quotation tag is used for quotations of more than one paragraph.
Following is the beginning of \emph{Alice's Adventures in Wonderland }by
Lewis Carroll:

\begin{quotation}
Alice was beginning to get very tired of sitting by her sister on the bank,
and of having nothing to do: once or twice she had peeped into the book her
sister was reading, but it had no pictures or conversations in it, 'and what
is the use of a book,' thought Alice 'without pictures or conversation?'

So she was considering in her own mind (as well as she could, for the hot
day made her feel very sleepy and stupid), whether the pleasure of making a
daisy-chain would be worth the trouble of getting up and picking the
daisies, when suddenly a White Rabbit with pink eyes ran close by her.

There was nothing so very remarkable in that; nor did Alice think it so very
much out of the way to hear the Rabbit say to itself, 'Oh dear! Oh dear! I
shall be late!' (when she thought it over afterwards, it occurred to her
that she ought to have wondered at this, but at the time it all seemed quite
natural); but when the Rabbit actually took a watch out of its
waistcoat-pocket, and looked at it, and then hurried on, Alice started to
her feet, for it flashed across her mind that she had never before seen a
rabbit with either a waistcoat-pocket, or a watch to take out of it, and
burning with curiosity, she ran across the field after it, and fortunately
was just in time to see it pop down a large rabbit-hole under the hedge.

In another moment down went Alice after it, never once considering how in
the world she was to get out again.
\end{quotation}

Use the Verbatim tag when you want \LaTeX{} to preserve spacing, perhaps
when including a fragment from a program such as:
\begin{verbatim}
#include <iostream>        // < > is used for standard libraries.
void main(void)            // "main" method always called first.
{
  cout << "Hello World.";  // Send to output stream.
}
\end{verbatim}

\subsection{Mathematics and Text}

Let $H$ be a Hilbert space, $C$ be a closed bounded convex subset of $H$, $T$
a nonexpansive self map of $C$. Suppose that as $n\rightarrow \infty $, $%
a_{n,k}\rightarrow 0$ for each $k$, and $\gamma _{n}=\sum_{k=0}^{\infty
}\left( a_{n,k+1}-a_{n,k}\right) ^{+}\rightarrow 0$. Then for each $x$ in $C$%
, $A_{n}x=\sum_{k=0}^{\infty }a_{n,k}T^{k}x$ converges weakly to a fixed
point of $T$ .

The numbered equation 
\begin{equation}
u_{tt}-\Delta u+u^{5}+u\left| u\right| ^{p-2}=0\text{ in }\mathbf{R}%
^{3}\times \left[ 0,\infty \right[  \label{eqn1}
\end{equation}
is automatically numbered as equation \ref{eqn1}.

\subsection{List Environments}

You can create numbered, bulleted, and description lists using the Item Tag
popup list on the Tag toolbar.

\begin{enumerate}
\item List item 1

\item List item 2

\begin{enumerate}
\item A list item under a list item.

The typeset style for this level is different than the screen style. The
screen shows a lower case alphabetic character followed by a period while
the typeset style uses a lower case alphabetic character surrounded by
parentheses.

\item Just another list item under a list item.

\begin{enumerate}
\item Third level list item under a list item.

\begin{enumerate}
\item Fourth and final level of list items allowed.
\end{enumerate}
\end{enumerate}
\end{enumerate}
\end{enumerate}

\begin{itemize}
\item Bullet item 1

\item Bullet item 2

\begin{itemize}
\item Second level bullet item.

\begin{itemize}
\item Third level bullet item.

\begin{itemize}
\item Fourth (and final) level bullet item.
\end{itemize}
\end{itemize}
\end{itemize}
\end{itemize}

\begin{description}
\item[Description List] Each description list item has a term followed by
the description of that term. Double click the term box to enter the term,
or to change it.

\item[Bunyip] Mythical beast of Australian Aboriginal legends.
\end{description}

\subsection{Theorem-like Environments}

The following theorem-like environments (in alphabetical order) are
available in this style.

\begin{acknowledgement}
This is an acknowledgement
\end{acknowledgement}

\begin{algorithm}
This is an algorithm
\end{algorithm}

\begin{axiom}
This is an axiom
\end{axiom}

\begin{case}
This is a case
\end{case}

\begin{claim}
This is a claim
\end{claim}

\begin{conclusion}
This is a conclusion
\end{conclusion}

\begin{condition}
This is a condition
\end{condition}

\begin{conjecture}
This is a conjecture
\end{conjecture}

\begin{corollary}
This is a corollary
\end{corollary}

\begin{criterion}
This is a criterion
\end{criterion}

\begin{definition}
This is a definition
\end{definition}

\begin{example}
This is an example
\end{example}

\begin{exercise}
This is an exercise
\end{exercise}

\begin{lemma}
This is a lemma
\end{lemma}

\begin{proof}
This is the proof of the lemma.
\end{proof}

\begin{notation}
This is notation
\end{notation}

\begin{problem}
This is a problem
\end{problem}

\begin{proposition}
This is a proposition
\end{proposition}

\begin{remark}
This is a remark
\end{remark}

\begin{solution}
This is a solution
\end{solution}

\begin{summary}
This is a summary
\end{summary}

\begin{theorem}
This is a theorem
\end{theorem}

\begin{proof}[Proof of the Main Theorem]
This is the proof.
\end{proof}

\subsection{Compiling a Bibliography}

\begin{enumerate}
\item Save this document using a name of your choosing.

\item Choose \textsf{Typeset, Compile.}

\item Check \textsf{Generate a Bibliography.}

\item Choose \textsf{OK.}

\item Choose \textsf{Typeset, Preview.}
\end{enumerate}

\subsection{Sample Citations and Bibliography}

BibTeX has been selected for the bibliography choice in this shell document.
One of the sample BibTeX databases included with SW has been selected and
some citations added in the next sentence. This sentence refers to the
article \cite{article-full}, the book \cite{book-full}, the booklet \cite%
{booklet-full}, the article \cite{incollection-full} in a collection, the
article \cite{inproceedings-full} in a proceedings, and to the thesis \cite%
{phdthesis-full}. The bibliography section for these citations comes next.

\bibliographystyle{achemso}
\bibliography{xampl}

\section{Introduction}

\label{achemsodoc}Author: Mats Dahlgren

Since the journals of The American Chemical Society (ACS) do not accept
manuscripts which are using the standard BibTeX formats, the need for such a
style package has been around for a while. I do sincerely hope that what I
have accomplished with this package is $(i)$~that using \LaTeX{} and BibTeX
should be as easy as any other manuscript preparation system when preparing
manuscripts for an ACS journal; $(ii)$~that the package is doing this
correctly. The reference handeling in the reference list and in the text is
aimed to be in accordance with the requirements of \textit{Accounts of
Chemical Research, Chemical Reviews, Inorganic Chemistry, Journal of the
American Chemical Society, Journal of Chemical Information and Computer
Sciences, Journal of Medicinal Chemistry, Journal of Organic Chemistry, The
Journal of Physical Chemistry, Langmuir, Macromolecules,} and \textit{%
Organometallics}. (All these journals use the same reference format.)

There is already a LaTeX 2.09 and BibTeX style package called \textsf{%
acsarticle} and \texttt{acs.bst}, which are not ``ACS'' as in `American
Chemical Society' (rather, this package is formatting the output according
to the instructions of \textit{Advances in Control Systems}). Hence, \textit{%
this} new package had to be given another name. The name of choice was then 
\textsf{achemso}, which is made from the words ``\textit{A}merican \textit{%
Chem}ical \textit{So}ciety''.

The \textsf{achemso} package has been tested with \LaTeX{} of 1997/12/01
patch level 1 running \TeX{} 3.14159 using MiK\TeX{} 1.09 containing BibTeX
version 0.99c, all under Windows 95. Please send bug reports (see below),
corrections, additions, suggestions, \textit{etc.}\ to me at \texttt{%
matsd@sssk.se}.

The present version of the package is~0.99{}, of 1997/02/22 . It differs
from previous versions in a change of the handeling of \texttt{phdthesis}, 
\texttt{mastersthesis}, and \texttt{book}. Version~0.1 produced an error for
these entry types if the optional field \texttt{institution} was empty. This
bug has been removed.\footnote{%
Thanks to Johan Fr\"{o}berg (\texttt{emgion@physchem.kth.se}) for bringing
my attention to the problem.} In the \texttt{book} entry type, there was a
comma which should be a semicolon. With version~0.3, Donald Arsenau's
package \textsf{overcite} is loaded, making collapsing and sorting of
references automatic. Version~0.4 was modified to work better with the
December 1995 release of \LaTeX{}; it also has the \texttt{.bst} files
incorporated in \texttt{achemso.dtx}. Version~1.0 was issued due to some new
pre-defined journal named and enhanced documentation.

This userguide is also available in \TEXTsymbol{\backslash}%
texttt\{.pdf\}-format on the internet. It is found from my \LaTeX{} web
page: 
%TCIMACRO{%
%\hyperref{http://www.homenet.se/matsd/latex/}{http://www.homenet.se/matsd/latex/}{}{http://www.homenet.se/matsd/latex/} }%
%BeginExpansion
\msihyperref{http://www.homenet.se/matsd/latex/}{http://www.homenet.se/matsd/latex/}{}{http://www.homenet.se/matsd/latex/}
%EndExpansion
(use Ctrl+Click to follow the link).

\subsection{Note}

This is not an official American Chemical Society package. Therefore, some
items in the reference list may not come out perfectly according to the ACS
style. There is no guarantee that the predefined journal abbreviations are
correct.

\section{The LaTeX Package}

To work properly, a call to the \textsf{achemso} package must be done in a 
\texttt{usepackge} command in the LaTeX input file. The \textsf{achemso}
package can take three options, namely \texttt{note}, \texttt{list}, and 
\texttt{number}. The \texttt{bibliographystyle} command required for a
document with a bibliography is included in the LaTeX package file. The
package options can be used together in any combination. Also, the \textsf{%
overcite} package must be installed on the system for \textsf{achemso} to
work properly.

The \textsf{achemso} package redefines the \texttt{bibliography} environment
to start a new page. If the \texttt{note} option is used, the heading of the
reference list will read ``References and Notes'' instead of ``References'',
which is the default. (If used together with the \textsf{babel} package,
there may be strange results. This is a minor problem since the ACS only
accepts manuscripts in English.)

In the text, references are printed as superscripts with numbers. If there
are more than two consecutive numbers at one occurence, the numbers will
appear ``collapsed'', that is like this$^{11-14}$ rather than.$%
^{11,12,13,14} $ The citations can be given in any order, thanks to the
sorting done by the \textsf{overcite} package, which is loaded by the 
\textsf{achemso} package. If only the reference number is desired, such as
in ``\ldots{} see reference 13.'', the command \texttt{citenum\{\}\textit{key%
}} can be used. \texttt{citenum} emits only the numeral of the reference
without any formatting.

\subsection{The \texttt{note} option}

\label{noteoption}

The option \texttt{note} is to be used when there are notes (of
footnote-type) in the text. The text of these notes have to be placed in a
BibTeX entry of the kind \texttt{remark}\footnote{%
The term ``\texttt{note}'' is already used for an entry field in BibTeX .}
(see also below). I recommend a seperate \texttt{.bib}-file for the notes,
which is then called for in the \texttt{bibliography}-command. In this way,
the notes will be sorted in automatically in the proper place in the
reference list, while they still are easily kept out of the listings of the
entire BibTeX data base. Also, the heading of the bibliography will read
``References and Notes'' and not only ``References''.

\subsection{The \texttt{list} option}

To produce a listing of the entire BibTeX data base, use the \texttt{list}
option for the \textsf{achemso} package. It makes use of another \texttt{.bst%
}-file, called \texttt{achemsol.bst}, which is also included in the \textsf{%
achemso} package distribution. It only differs from \texttt{achemso.bst} in
that it adds the information of the field \texttt{annotate} to every item in
the list when present. \texttt{annotate} is an optional field for all entry
types.

When usnig the \texttt{list} option, the BibTeX keys will be typesetted in a
framed box in the left margin. This feature is adapted from \texttt{%
showtags.sty} by Nelson H.\ F.\ Beebe (\texttt{beebe@math.utah.edu}).%
\footnote{%
Many thanks to Nelson H.\ F.\ Beebe for letting me incorporate his code into 
\textsf{achemso}.}

For use in the \texttt{annotate} field is a LaTeX macro \texttt{refin}
defined by \texttt{achemso.bst} and \texttt{achemsol.bst}. This command
takes one argument (normally text) which is preceeded by the text ``\textbf{%
Referenced in: }''. \texttt{refin} is meant to facilitate in keeping track
of in which publication(s) the bibliography item has been cited. Personally,
I keep information in my data-base on in which of my own publications I have
cited the article (or whatever) in question. It is for this purpouse the 
\texttt{refin} command exsists.

\subsection{The \texttt{number} option}

If a section number is desired for the bibliography in a document, the 
\texttt{number} option should be used with the \textsf{achemso} package. The
only difference with respect to the default version is the adding of a
section number to the heading. (It makes use of \texttt{section} rather than 
\texttt{section*} as the default.)

\section{The BibTeX listing file}

To obtain a listing of a BibTeX data base in an easy way, only very minor
changes are needed to the file \texttt{acslist.tex}.\footnote{%
You must specify the name(s) of your \texttt{.bib}-file(s).} When running
LaTeX on this file, you will be provided with a complete listing of the data
base as the output, with the labels used in framed boxes to the left of the
data base entries. \texttt{acslist.tex} uses the \texttt{list} option of the
package file (surprise!). To get rid of the sectioning heading in the data
base listing, you can include the following line in \texttt{acslist.tex}:%
\newline
\texttt{\ renewcommand\{refname\}\{ \}}\newline

\section{The BibTeX style}

The \texttt{.bst}-files in the \textsf{achemso} package are based on the
standard BibTeX style \texttt{unsrt.bst}.

\subsection{Extra Entry Types}

The BibTeX style \texttt{achemso} defines three non-standard entry types, 
\texttt{inpress}, \texttt{submitted}, and \texttt{remark}. The \texttt{%
inpress} and \texttt{submitted} entries are intended for journal articles
which have not yet been published, but are ``in press'' or ``submitted'',
respectively. The \texttt{remark} entry type is to be used to obtain notes
intermixed in the reference list (see also section~\ref{noteoption} above).

\subsection{The Data Base Entries}

This section explains which data base entry types are recognized by \texttt{%
achemso} and which fields are required, optional, and ignored for the
different entries. Note that \texttt{achemso} is not as robust as many other
BibTeX styles, and may in some cases not warn that a required field is
missing. The only way to detect this is to look in the final bibliography
listing, which for the incomplete entries is likely to look strange. The
following entries are recognized:\\[\baselineskip]
\hfill 
\begin{tabular}{|ll|}
\hline
\textbf{Entry type} & \textbf{Entry fields} \\ \hline
\texttt{article} & \textsf{R:}~\texttt{author}, \texttt{journal}, \texttt{%
year}, \texttt{volume}, \texttt{pages} \\ 
& \textsf{O:}~\texttt{note} \\ \hline
\texttt{book} & \textsf{R:}~\texttt{author} or \texttt{editor}, \texttt{title%
} or \texttt{booktitle}, \texttt{publisher}, \texttt{address}, \texttt{year}
\\ 
& \textsf{O:}~\texttt{volume}, \texttt{series}, \texttt{edition}, \texttt{%
note} \\ \hline
\texttt{booklet} & \textsf{R:}~\texttt{author}, \texttt{title}, \texttt{year}
\\ 
& \textsf{O:}~\texttt{howpublished}, \texttt{address}, \texttt{note} \\ 
\hline
\texttt{conference} & \textsf{R:}~\texttt{author}, \texttt{title}, \texttt{%
booktitle}, \texttt{editor}, \texttt{publisher}, \texttt{address}, \texttt{%
year} \\ 
& \textsf{O:}~\texttt{volume}, \texttt{edition}, \texttt{type}, \texttt{%
chapter}, \texttt{pages}, \texttt{note} \\ \hline
\texttt{inbook} & \textsf{R:}~\texttt{author}, \texttt{title}, \texttt{%
booktitle}, \texttt{editor}, \texttt{publisher}, \texttt{address}, \texttt{%
year} \\ 
& \textsf{O:}~\texttt{volume}, \texttt{edition}, \texttt{type}, \texttt{%
chapter}, \texttt{pages}, \texttt{note} \\ \hline
\texttt{incollection} & \textsf{R:}~\texttt{author}, \texttt{title}, \texttt{%
booktitle}, \texttt{editor}, \texttt{publisher}, \texttt{address}, \texttt{%
year} \\ 
& \textsf{O:}~\texttt{volume}, \texttt{edition}, \texttt{note} \\ \hline
\texttt{inpress} & \textsf{R:}~\texttt{author}, \texttt{journal} \\ 
& \textsf{O:}~\texttt{note} \\ \hline
\texttt{inproceedings} & \textsf{R:}~\texttt{author}, \texttt{title}, 
\texttt{booktitle}, \texttt{editor}, \texttt{publisher}, \texttt{address}, 
\texttt{year} \\ 
& \textsf{O:}~\texttt{volume}, \texttt{edition}, \texttt{type}, \texttt{%
chapter}, \texttt{pages}, \texttt{note} \\ \hline
\texttt{manual} & \textsf{R:}~\texttt{title} \\ 
& \textsf{O:}~\texttt{author}, \texttt{organization}, \texttt{address}, 
\texttt{edition}, \texttt{year}, \texttt{note} \\ \hline
\texttt{mastersthesis} & \textsf{R:}~\texttt{author}, \texttt{title}, 
\texttt{school}, \texttt{year} \\ 
& \textsf{O:}~\texttt{institution}, \texttt{note} \\ \hline
\texttt{misc} & \textsf{R:}~none \\ 
& \textsf{O:}~\texttt{author}, \texttt{title}, \texttt{howpublished}, 
\texttt{year}, \texttt{note} \\ \hline
\texttt{phdthesis} & \textsf{R:}~\texttt{author}, \texttt{title}, \texttt{%
school}, \texttt{year} \\ 
& \textsf{O:}~\texttt{type}, \texttt{institution}, \texttt{note} \\ \hline
\texttt{proceedings} & \textsf{R:}~\texttt{editor} or \texttt{organization}, 
\texttt{title} or \texttt{booktitle}, \texttt{year} \\ 
& \textsf{O:}~\texttt{volume}, \texttt{edition}, \texttt{publisher}, \texttt{%
address}, \texttt{note} \\ \hline
\texttt{remark} & \textsf{R:}~\texttt{note} \\ 
& \textsf{O:}~none \\ \hline
\texttt{submitted} & \textsf{R:}~\texttt{author}, \texttt{journal} \\ 
& \textsf{O:}~\texttt{note} \\ \hline
\texttt{techreport} & \textsf{R:}~\texttt{author}, \texttt{title}, \texttt{%
institution} \\ 
& \textsf{O:}~\texttt{type}, \texttt{number}, \texttt{address}, \texttt{note}
\\ \hline
\texttt{unpublished} & \textsf{R:}~\texttt{author } \\ 
& \textsf{O:}~\texttt{year}, \texttt{note} \\ \hline
\end{tabular}
\hfill\newline
\textsf{R:}~denotes required fields; \textsf{O:}~denotes optional filelds.\\[%
\baselineskip]
The following fields are avaliable: \texttt{address}, \texttt{author}, 
\texttt{booktitle}, \texttt{chapter}, \texttt{edition}, \texttt{editor}, 
\texttt{howpublished}, \texttt{institution}, \texttt{journal}, \texttt{key}, 
\texttt{month}, \texttt{note}, \texttt{number}, \texttt{organization}, 
\texttt{pages}, \texttt{publisher}, \texttt{school}, \texttt{series}, 
\texttt{title}, \texttt{type}, \texttt{volume}, \texttt{year}, and \texttt{%
annotate}. All the fields have their general meanings. All fields not
specified above are ignored for that entry type.

\subsection{Predefined Journal Abbriviations}

In the \textsf{achemso} package, some journal names are defined. These
journals are the journals of ACS, some other frequently cited in Physical
Chemsitry, and the ten most cited according to \textit{Chemical Abstracs}.
The predefined ACS journal names are:\\[\baselineskip]
\begin{tabular}{|lll|}
\hline
Abbreviation & Appearance in list & Full journal name \\ \hline
\texttt{acchemr} & \textit{Acc.\ Chem.\ Res.} & Accounts of Chemical Research
\\ 
\texttt{aacsa} & \textit{Adv.\ {ACS} Abstr.} & Advance ACS Abstracts \\ 
\texttt{anchem} & \textit{Anal.\ Chem.} & Analytical Chemistry \\ 
\texttt{bioch} & \textit{Biochemistry} & Biochemistry \\ 
\texttt{bicoc} & \textit{Bioconj.\ Chem.} & (?) Bioconjugate Chemistry \\ 
\texttt{bitech} & \textit{Biotechnol.\ Progr.} & (?) Biotechnology Progress
\\ 
\texttt{chemeng} & \textit{Chem.\ Eng.\ News} & Chemical \&{} Engineering
News \\ 
\texttt{chs} & \textit{Chem.\ Health Safety} & (?) Chemical Health \&{}
Safety \\ 
\texttt{crt} & \textit{Chem.\ Res.\ Tox.} & (?) Chemical Research in
toxicology \\ 
\texttt{chemrev} & \textit{Chem.\ Rev.} & Chemical Reviews \\ 
\texttt{cmat} & \textit{Chem.\ Mater.} & Chemistry of Materials \\ 
\texttt{chemtech} & \textit{CHEMTECH} & CHEMTECH \\ 
\texttt{enfu} & \textit{Energy \&{} Fuels} & (?) Energy \&{} Fuels \\ 
\texttt{envst} & \textit{Environ.\ Sci.\ Technol.} & Environmental Science
and Technology \\ 
\texttt{iecf} & \textit{Ind.\ Eng.\ Chem.\ Fundam.} & Industrial \&{}
Engineering Chemistry \\ 
&  & \hfill Fundamentals \\ 
\texttt{iecpdd} & \textit{Ind.\ Eng.\ Chem.\ Proc.} & Industrial \&{}
Engineering Chemistry \\ 
& \hfill\textit{Des.\ Dev.} & \hfill Process Design and Development \\ 
\texttt{iecprd} & \textit{Ind.\ Eng.\ Chem.\ Prod.} & Industrial \&{}
Engineering Chemistry \\ 
& \hfill\textit{Res.\ Dev.} & \hfill Product Research and Development \\ 
\texttt{iecr} & \textit{Ind.\ Eng.\ Chem.\ Res.} & Industrial \&{}
Engineering Chemistry \\ 
&  & \hfill Research \\ 
\texttt{inor} & \textit{Inorg.\ Chem.} & Inorganic Chemistry \\ 
\texttt{jafc} & \textit{J.~Agric.\ Food Chem.} & Journal of Agricultural and
Food Chemistry \\ 
\texttt{jacs} & \textit{J.~Am.\ Chem.\ Soc.} & Journal of the American
Chemical Society \\ 
\texttt{jced} & \textit{J.~Chem.\ Eng.\ Data} & Journal of Chemical and
Engineering Data \\ 
\texttt{jcics} & \textit{J.~Chem.\ Inf.\ Comput.\ Sci.} & Journal of
Chemical Information and \\ 
&  & \hfill Computer Sciences \\ 
\texttt{jmc} & \textit{J.~Med.\ Chem.} & Journal of Medicinal Chemistry \\ 
\texttt{joc} & \textit{J.~Org.\ Chem.} & Journal of Organic Chemistry \\ 
\texttt{jps} & \textit{J.~Pharm.\ Sci.} & Journal of Pharmaceutical Sciences
\\ 
\texttt{jpcrd} & \textit{J.~Phys.\ Chem.\ Ref.\ Data} & Journal of Physical
and Chemical \\ 
&  & \hfill Reference Data \\ 
\texttt{jpc} & \textit{J.~Phys.\ Chem.} & The Journal of Physical Chemistry
\\ 
\texttt{lang} & \textit{Langmuir} & Langmuir \\ 
\texttt{macro} & \textit{Macromolecules} & Macromolecules \\ 
\texttt{orgmet} & \textit{Organometallics} & Organometallics \\ \hline
\end{tabular}

\noindent Other predefined journal names are:\newline
\begin{tabular}{|lll|}
\hline
Abbreviation & Appearance in list & Full journal name \\ \hline
\texttt{jft} & \textit{J.~Chem.\ Soc., Faraday Trans.} & Journal of the
Chemical Society, \\ 
&  & \hfill Faraday Transactions \\ 
\texttt{jft1} & \textit{J.~Chem.\ Soc., Faraday Trans. 1} & Journal of the
Chemical Society, \\ 
&  & \hfill Faraday Transactions 1 \\ 
\texttt{jft2} & \textit{J.~Chem.\ Soc., Faraday Trans. 2} & Journal of the
Chemical Society, \\ 
&  & \hfill Faraday Transactions 2 \\ 
\texttt{tfs} & \textit{Trans.\ Faraday Soc.} & Transactions of the Faraday
Society \\ 
\texttt{jcis} & \textit{J.~Colloid Interface Sci.} & Journal of Colloid and
Interface Science \\ 
\texttt{acis} & \textit{Adv.~Colloid Interface Sci.} & Advances in Colloid
and Interface Science \\ 
\texttt{cs} & \textit{Colloids Surf.} & Colloid and Surfaces \\ 
\texttt{csa} & \textit{Colloids Surf.\ A:\ Physiochem.} & Colloid and
Surfaces A:\ Physiochemical \\ 
& \hfill \textit{Eng.\ Aspects} & and Engineering Aspects \\ 
\texttt{csb} & \textit{Colloids Surf.\ B:\ Biointerfaces} & Colloid and
Surfaces B:\ Biointerfaces \\ 
\texttt{pcps} & \textit{Progr.\ Colloid Polym.\ Sci.} & Progress in Colloid
and Polymer Science \\ 
\texttt{jmr} & \textit{J.~Magn.\ Reson.} & Journal of Magnetic Resonance \\ 
\texttt{jmra} & \textit{J.~Magn.\ Reson.\ A} & Journal of Magnetic
Resonance, Series A \\ 
\texttt{jmrb} & \textit{J.~Magn.\ Reson.\ B} & Journal of Magnetic
Resonance, Series B \\ 
\texttt{sci} & \textit{Science} & Science \\ 
\texttt{nat} & \textit{Nature (London)} & Nature \\ 
\texttt{jcch} & \textit{J.~Comput.\ Chem.} & Journal of Computational
Chemistry \\ 
\texttt{cca} & \textit{Croat.\ Chem.\ Acta} & Croatica Chemica Acta \\ 
\texttt{poly} & \textit{Polymer} & Polymer \\ 
\texttt{ajp} & \textit{Am.\ J.\ Phys.} & American Journal of Physics \\ 
\texttt{rsi} & \textit{Rev.\ Sci.\ Instrum.} & Review of Scientific
Instruments \\ 
\texttt{jcp} & \textit{J.~Chem.\ Phys.} & Journal of Chemical Physics \\ 
\texttt{cpl} & \textit{Chem.\ Phys.\ Lett.} & Chemical Physics Letters \\ 
\texttt{molph} & \textit{Mol.\ Phys.} & Molecular Physics \\ 
\texttt{pac} & \textit{Pure Appl.\ Chem.} & Pure and Applied Chemistry \\ 
\texttt{jbc} & \textit{J.~Biol.\ Chem.} & Journal of Biological Chemistry \\ 
\texttt{tl} & \textit{Tetrahedron Lett.} & Tetrahedron Letters \\ 
\texttt{psisoe} & \textit{Proc.\ SPIE-Int.\ Soc.\ Opt.\ Eng.} & (?) \\ 
\texttt{prb} & \textit{Phys.\ Rev.\ B:\ Condens.\ Matter} & Physics Review
B:\ Condensed Matter \\ 
\texttt{jap} & \textit{J.~Appl.\ Phys.} & Journal of Applied Physics \\ 
\texttt{pnac} & \textit{Proc.\ Natl.\ Acad.\ Sci.\ U.~S.~A.} & Proceedings
of the National Academy \\ 
&  & \hfill of Sciences of the U.~S.~A. \\ 
\texttt{bba} & \textit{Biochim.\ Biophys.\ Acta} & Biochimica et Boiphysica
Acta \\ 
\texttt{nar} & \textit{Nucleic Acid Res.} & Nucleic Acid Research (?) \\ 
\hline
\end{tabular}
\newpage

\section{Known Problems and Limitations}

At the present (1997/02/22), the following problems are known:

\begin{itemize}
\item When printing a data base listing using the \texttt{list} option (such
as when type-setting \texttt{acslist.tex}), there may on rare occations
occur a page-break between the framed box of the label and the BibTeX entry
itself.

\item Some of the pre-defined journal abbreviations may not be entirely
correct, I have not been able to check some of them (since those journals
are not avaliable at our library). If someone finds an error in the journal
names, I'll be most grateful to be informed on how the correct abbreviation
is.

\item Compatibility with other LaTeX packages, both standard and
non-standard, may be poor.
\end{itemize}

\end{document}
